\setauthor{\firstauthor}

\section{Ausgangslage}
Die Vorbereitung auf die Matura stellt für viele Schüler:innen eine besondere Herausforderung dar.
Der Überblick über verschiedene Fächer, Themenpools und Unterlagen geht schnell verloren, 
während der Stresspegel vor den Prüfungen steigt. Besonders schwierig gestaltet sich dabei 
die Strukturierung des umfangreichen Lernstoffs sowie die Suche nach geeigneten Lernmaterialien.\par
Aktuell werden Mitschriften und Lernmaterialien an der HTL Leonding über verschiedene Kanäle verteilt - 
über Microsoft Teams, Moodle, Schul-Laufwerke oder in Papierform.
Dies führt dazu, dass wertvolle Materialien schwer auffindbar sind oder mit der Zeit verloren gehen. 
Schüler:innen müssen Unterlagen aus verschiedenen Jahren mühsam zusammensuchen und sortieren.\par
Zusätzlich kommt hinzu, dass es für HTL-spezifische Fächer online nur begrenzte Lernhilfen gibt. 
Anders als bei allgemeinbildenden Fächern finden Schüler:innen einer HTL oft keine passenden externen Ressourcen.\par
Um diesen Problemen entgegenzuwirken, entwickeln zwei Diplomarbeitsteams gemeinsam die Web-Plattform 
\textit{LearnHub}, die es Schüler:innen und Lehrer:innen der HTL ermöglicht, Mitschriften zentral zu 
verwalten, um sich gezielt, ohne zusätzlichen Stress, auf die Matura vorzubereiten.

Die Entwicklung von LearnHub wurde in Zusammenarbeit mit einem zweiten Diplomarbeitsteam der HTL Leonding realisiert. 
Während sich die vorliegende Arbeit auf das Content-Management-System (Maturaupload und Fragenkonfigurator) konzentriert, 
entwickelte das Partnerteam den Maturatrainer - ein Lern- und Prüfungsmodus.

\section{Marktanalyse}
Es befinden sich bereits verschiedene Lernplattformen auf dem Markt. Im Folgenden werden vier dieser Konkurrenten genauer analysiert, 
um die Besonderheiten von LearnHub hervorzuheben.\par
Die meisten dieser Plattformen sind kostenpflichtig oder bieten in ihrer kostenlosen Version nur eingeschränkte Funktionen. 
Beispielhaft sind hier einige Konkurrenzprodukte aufgeführt: 
\subsubsection{School Fox, Study Smarter, Evernote, Studyflix}

\subsection{Schlussfolgerung}
Der Vergleich mit verschiedenen Lernplattformen zeigt, dass zwar viele Tools für die Lernunterstützung existieren, 
jedoch keine Plattform die speziell auf die Bedürfnisse der Schüler:innen der HTL Leonding zugeschnitten ist.\par
Besonders fehlt bei analysierten Plattformen:
\begin{itemize}
    \item Schultypspezifische Mitschriften
    \item Strukturierte Verwaltung nach Fächern und Themenpool
    \item Integriertem Fragenkonfigurator
    \item Kostenfreier Zugang 
\end{itemize}

LearnHub schließt diese Lücke, indem es eine Lösung für die HTL Leonding bietet. Die Lernplattform bietet  
das erstellen von Maturaunterlagen und Maturafragen von HTL-spezifischen Maturafächern. 

\section{Motivation und Zielsetzung}
LearnHub soll Schüler:innen der HTL Leonding dabei helfen, sich strukturiert und effizient auf die Matura vorzubereiten, 
indem sie Zugriff auf Lernmaterialien zum Lernen für die Matura erhalten.\par
Im Rahmen dieser Diplomarbeit wird das \textbf{Content-Management-System} entwickelt, welches zwei zentrale Punkte beinhaltet:

\subsubsection{Maturaupload}
Der Maturaupload ermöglicht die strukturierte Verwaltung von Mitschriften und Lernunterlagen:
\begin{enumerate}
    \item \textbf{Upload und Verwaltung}: Schüler:innen können eigene Mitschriften hochladen, bearbeiten und löschen.
        Lehrer:innen haben erweiterte Rechte und können alle Mitschriften verwalten.
    \item \textbf{Freigabeprozess}: Hochgeladene Mitschriften müssen von Lehrer:innen freigegeben werden, bevor sie für alle sichtbar sind.
    \item \textbf{Strukturierung}: Mitschriften werden nach Fächern und Themenpools organisiert, sodass Schüler:innen schnell die 
        benötigten Unterlagen finden können. Nur Lehrer:innen können neue Fächer und Themenpools anlegen und verwalten.
    \item \textbf{Nutzung}: Alle freigegebenen Mitschriften können von Schüler:innen eingesehen werden.
\end{enumerate}

\subsubsection{Fragenkonfigurator}
Der Fragenkonfigurator bietet ein flexibles System zum Erstellen und Lernen von Fragen:
\begin{enumerate}
    \item \textbf{Private Fragen}: Schüler:innen können eigene private Fragen erstellen, bearbeiten und löschen.
        Diese sind nur für die/den Nutzer:in selbst sichtbar. Diese Fragen dienen zur persönlichen Vorbereitung.
    \item \textbf{Öffentliche Fragen}: Lehrer:innen können öffentliche Fragen erstellen, bearbeiten und löschen. 
        Diese sind für alle Nutzer:innen sichtbar.
    \item \textbf{Lernmodus}: Schüler:innen können sowohl mit öffentlichen als auch mit ihren eigenen privaten Fragen üben.
        Der Lernmodus funktioniert wie ein digitales Karteikartensystem, dabei können Fragen geübt werden, 
        die richtigen Antworten können eingeblendet werden.
\end{enumerate}
